\chapter*{Resumen}
\addcontentsline{toc}{chapter}{Resumen}

El crecimiento exponencial del espacio de Hilbert con el número de constituyentes cuánticos hace que numerosas tareas de información cuántica sean intratables mediante métodos clásicos exactos, lo que motiva el desarrollo de aproximaciones escalables. Esta tesis presenta dos enfoques de physics-informed deep learning que abordan este tipo de problemas incorporando los principios físicos subyacentes directamente en la arquitectura de la red neuronal y en la función objetivo de entrenamiento, garantizando que las soluciones aprendidas permanezcan consistentes con las leyes de la mecánica cuántica.

El primer enfoque estudia la dinámica de sistemas cuánticos cerrados gobernados por Hamiltonianos dependientes del tiempo, cuya evolución está descrita por la ecuación de von Neumann. En este contexto, el operador de evolución temporal se expresa mediante expansiones integrales infinitas. A diferencia del caso independiente del tiempo, donde la solución admite una forma exponencial simple, obtener una expresión análoga requiere la introducción del operador de ordenamiento temporal. En consecuencia, cuando el Hamiltoniano no conmuta consigo mismo en distintos instantes de tiempo, la complejidad matemática del problema aumenta considerablemente. El método propuesto reduce el costo computacional asociado a la simulación de estas dinámicas. A partir del operador de evolución temporal conocido únicamente en un conjunto disperso de instantes de tiempo, una Physics-Informed Neural Network interpola su evolución sobre todo el dominio temporal sin requerir conocimiento explícito del Hamiltoniano. En su lugar, el modelo incorpora la unitariedad como una restricción física y, para sistemas de mayor tamaño, predice un Hamiltoniano efectivo cuya evolución se propaga mediante integradores que preservan la estructura del sistema. Para sistemas de entre 2 y 8 qubits, el enfoque alcanza fidelidades promedio cercanas a 0.99, reproduce con precisión la dinámica incluso bajo un muestreo temporal poco denso y logra acelerar la inferencia aproximadamente en un orden de magnitud con respecto a un método de interpolación geodésica.

El segundo enfoque aborda el Quantum Marginal Problem, cuyo objetivo es reconstruir un estado cuántico global a partir de un conjunto dado de matrices densidad reducidas (marginales cuánticos), garantizando simultáneamente que todas ellas sean mutuamente compatibles. Este problema se vuelve cada vez más desafiante a medida que aumenta el tamaño del sistema debido al crecimiento exponencial del espacio de Hilbert y a las complejas restricciones de compatibilidad entre las marginales. Mediante la combinación de un convolutional denoising autoencoder con el Marginal Imposition Operator, y representando las matrices densidad como imágenes de dos canales junto con una función de pérdida que impone Hermiticidad, positividad y normalización, el modelo reconstruye estados globales físicamente válidos para sistemas de entre 3 y 8 qubits. Un esquema híbrido de reconstrucción recupera, en numerosos casos, estados cuyas marginales coinciden con las marginales objetivo con fidelidad perfecta, mientras que transfer learning permite extender el modelo a distintos tamaños de sistema con un reentrenamiento mínimo. Una vez entrenado, el enfoque realiza inferencias sustancialmente más rápidas que los solucionadores de programación semidefinida más avanzados, logrando además reconstruir estados en regímenes donde dichos métodos se vuelven computacionalmente impracticables.

En conjunto, estos resultados demuestran que la incorporación explícita de principios físicos dentro del proceso de aprendizaje permite que el deep learning actúe no como un reemplazo del modelamiento físico, sino como una extensión natural de éste. Al combinar aprendizaje basado en datos con la estructura matemática de la mecánica cuántica, los enfoques propuestos permiten abordar de manera eficiente problemas cuyo tratamiento clásico exacto resulta computacionalmente prohibitivo.